% Page budget: 0.35 pages | Owns: problem, method, principal evidence, conclusion, and index terms.
\begin{abstract}
\begin{itemize}
\item Problem: improve settling prediction for a position-dependent dual gantry without surrendering the meaning of its physical parameters
\item Method: combine a quasi-LPV first-principles baseline, dynamic parallel LFR augmentation, closed-loop rollout training, and state-level orthogonal projection
\item Evidence: compare baseline, unconstrained augmentation, orthogonal variants, and a black-box model on held-out positions and motion profiles
\item Main result: report predictive improvement together with the measured limits from short-horizon training, inactive added states, and non-orthogonal true missing dynamics
\item Conclusion: state when the proposed decomposition is useful and which conditions prevent interpreting it as true parameter recovery
\end{itemize}
\todo{Write the numerical abstract only after the final simulation and Telica result tables are frozen.}
\end{abstract}

\begin{IEEEkeywords}
grey-box identification, LPV systems, model augmentation, orthogonal projection, dual gantry
\end{IEEEkeywords}
